% ---------------------------------------------------------------------------
% Draft section: Compile-Time Enforcement of Register Access Constraints
% Pair of docs/paper_enforcement_section_draft.md (this .tex is the paper
% source; the .md is the readable review copy).
%
% Requires in the preamble:
%   \usepackage{booktabs}
%   \usepackage{listings}
%   \usepackage{xcolor}   % optional, for the listing frame color
% Suggested listing setup (once, in the preamble):
%   \lstset{basicstyle=\ttfamily\footnotesize, breaklines=true,
%           frame=single, framerule=0.3pt, columns=fullflexible,
%           keepspaces=true, showstringspaces=false}
%
% Numbers are reproducible: grammar/corpus figures from
% docs/constraints_corpus_stats.md and docs/quote_anchoring_stats.md;
% enforcement and HAL figures pinned by applications/pac_codegen/
% test_codegen.py (test_illegal_programs_rejected, test_hal_demo).
% ---------------------------------------------------------------------------

\section{Compile-Time Enforcement of Register Access Constraints}
\label{sec:enforcement}

STMicroelectronics' reference manuals for the STM32 family state, in the
description of the I\textsuperscript{2}C peripheral's first control register:

\begin{quote}
``When the STOP, START or PEC bit is set, the software must not perform any
write access to I2C\_CR1 before this bit is cleared by hardware. This is to
avoid setting a second STOP, START or PEC request.''
\end{quote}

This is an \emph{access-ordering constraint}: a rule about \emph{when} a
register may be touched, not about \emph{what} the register contains. SVD
files---the machine-readable register descriptions from which peripheral
access crates (PACs) are generated---have no vocabulary for such rules, so
today they live only in prose and in driver authors' heads. Our extraction
pipeline recovers this sentence and its meaning from the manual; the four
subsections below follow it from structured form to a compile error in an
unmodified, widely-used driver library.

\subsection{The Register Access Constraint Grammar}
\label{sec:enforcement:grammar}

The grammar is the contract between the extraction LLM and everything
downstream: a small, closed target language that is easy for a model to emit,
mechanically validatable without any AI in the loop, and compilable into
enforcement. Three design rules shape it.

\paragraph{Closed vocabularies, structured values.}
Every enumerable field is a fixed choice set (operation $\in$ \{read, write,
modify, any\}; severity $\in$ \{error, warning\}), and numeric values are
lists of integers validated against the SVD's field widths---never free
text. An earlier iteration allowed strings such as
\texttt{equals:0b01|0b10|0b11}; free text of this kind both invites drift
and, when spliced into generated code, silently changes meaning under the
host language's operator precedence. In the current grammar such a value is
the list $[1, 2, 3]$.

\paragraph{A quote anchor.}
Every constraint carries the verbatim sentence it was extracted from.
Because the quote travels with the constraint, its authenticity is checkable
by plain string search against the manual---no retrieval, no judgment---and
the surrounding paragraph can be recovered as trusted context for later
validation.

\paragraph{An evidence dichotomy.}
Each condition declares \emph{who establishes the state}: \texttt{hardware}
(software may only observe it---the running example, where hardware clears
STOP) or \texttt{software} (the driver itself must establish it, e.g.\ ``this
register may be written only while the peripheral is disabled''). The
distinction cannot be inferred from the SVD; it exists only in prose, so it
must be captured at extraction time. As \S\ref{sec:enforcement:witnesses}
shows, it decides the \emph{shape} of the generated enforcement.

Constraints are a tagged union of eight kinds---state gates (the workhorse),
multi-step sequences (unlock keys), write-once locks, delays, read
side-effects, peripheral clock gates, value relations, and an explicit
\texttt{other} escape valve so a real requirement that fits no kind is
preserved verbatim rather than force-fitted or dropped. Enforceability is
\emph{computed} from the kind and evidence, never emitted by the model.
The running example in grammar form:

\begin{lstlisting}
{ "kind": "state_gate",
  "target_register": "CR1", "target_operation": "write",
  "preconditions": [
    {"register": "I2C_CR1", "field": "STOP",  "state": "cleared", "evidence": "hardware"},
    {"register": "I2C_CR1", "field": "START", "state": "cleared", "evidence": "hardware"},
    {"register": "I2C_CR1", "field": "PEC",   "state": "cleared", "evidence": "hardware"}],
  "postconditions": [],
  "severity": "error",
  "datasheet_text": "When the STOP, START or PEC bit is set, ..." }
\end{lstlisting}

\paragraph{Results.}
Applied to 30 STM reference manuals, the pipeline extracted 4{,}927 raw
constraints (4{,}362 after removing exact duplicates). Deterministic
validation---vocabulary, SVD name and width resolution, duplicate and
flag-semantics filters---accepts \textbf{2{,}857 enforceable state gates} and
rejects the remainder with machine-readable reasons (35.6\% under a
one-device-per-manual projection; 28.5\% when resolving against all of a
manual's devices---the dominant reject is a constraint naming a peripheral
the projected device does not have). Independently, \textbf{91.1\% of all
extracted constraints carry a quote genuinely present in the manual} (78.2\%
verbatim, 12.9\% near-verbatim); the unanchored 8.9\% are stitched or
paraphrased quotes and are routed to human review.

\subsection{Translating Constraints to Witness Tokens}
\label{sec:enforcement:witnesses}

The compiler cannot observe hardware, so it cannot enforce ``STOP is
clear.'' What it \emph{can} enforce is an ordering: \textbf{the
datasheet-prescribed check must happen before the constrained operation}.
The translation makes that ordering a type-system fact using \emph{witness
tokens}: zero-sized values that can only be created by performing the
prescribed check, and that the constrained operation consumes. For the
running example the generator emits, into the PAC's I\textsuperscript{2}C
module:

\begin{lstlisting}
/// State witness authorizing one write of CR1.
pub struct Cr1WriteWitness { _priv: () }   // cannot be forged: private field

impl Reg<CR1rs> {
    /// Read CR1 once and check every write precondition.
    pub fn check_write_ready(&self)
        -> Result<Cr1WriteWitness, Cr1ConstraintError>
    {
        let r = self.read();                // ONE fresh hardware read
        if !r.stop().bit_is_clear()  { return Err(StopNotCleared); }
        if !r.start().bit_is_clear() { return Err(StartNotCleared); }
        if !r.pec().bit_is_clear()   { return Err(PecNotCleared); }
        Ok(Cr1WriteWitness { _priv: () })
    }
    /// Check and write in one call -- the witness never escapes user code.
    pub fn write_when_ready<F>(&self, f: F)
        -> Result<u16, Cr1ConstraintError> { ... }
}
\end{lstlisting}

Four properties carry the guarantees. The witness is \emph{unforgeable} (its
only constructor is the check). It is \emph{affine}---consumed by the write,
so one check authorizes one operation; reusing it is a compile error. It is
\emph{operation-specific}: a write witness does not authorize a
read-modify-write, which has its own witness type. And the recommended entry
point (\texttt{write\_when\_ready}) mints and spends the witness inside one
call, so the check-to-use window is fixed by the library, not by caller
discipline.

Honesty about semantics: a witness is a \emph{witness}, not a proof---it
attests that the preconditions were observed true in one fresh read.
Time-of-check to time-of-use is inherited from the hardware contract itself:
memory-mapped I/O has no atomic check-and-act, so the manual's own prescribed
procedure carries the identical window, and for hardware-cleared flags like
STOP the race is benign in the safe direction (hardware only clears it). The
evidence dichotomy of \S\ref{sec:enforcement:grammar} selects the witness's
origin: \texttt{hardware} evidence produces a runtime check as above (a
\emph{state witness}); \texttt{software} evidence produces a token minted by
performing the required setup action (an \emph{action witness}), with no
runtime check at all.

\paragraph{Results.}
Of the 2{,}857 accepted gates, 2{,}243 translate to runtime-checked state
witnesses and 614 remain documentation-only (they carry no checkable
condition in the legacy corpus encoding); action witnesses appear once
re-extraction emits the evidence field, making that split a direct measure
of the re-extraction step's value.

\subsection{Applying the Changes to PAC Crates}
\label{sec:enforcement:pac}

\paragraph{How a PAC crate works today.}
A peripheral access crate is machine-generated (by the community tool
\texttt{svd2rust}) from the vendor's SVD file, one crate per family with a
build feature per device. Its safety model rests on two ideas. First,
\emph{ownership}: the crate exposes each peripheral as a singleton object
(\texttt{Peripherals::take()} yields it exactly once), so two parts of a
program cannot race on the same hardware by accident. Second, \emph{typed
registers}: every register is a value of type \texttt{Reg<REG>}, where
\texttt{REG} is a zero-sized marker type identifying that specific register,
and a small set of \textbf{capability traits} on the marker declares what the
hardware allows---\texttt{RegisterSpec} (its raw width), \texttt{Readable},
\texttt{Writable}, and \texttt{Resettable} (its documented reset value). The
register methods are defined once, generically, but each is available only
where the matching capability holds: \texttt{read()} returns a reader with
typed per-field accessors (\texttt{r.stop().bit\_is\_clear()});
\texttt{write(|w| \ldots)} composes a full register value starting from the
reset value through field proxies (\texttt{w.pe().enabled()}) that make
invalid bit patterns unrepresentable---or \texttt{unsafe} where the SVD
cannot vouch for them; \texttt{modify(|r, w| \ldots)} performs a
read-modify-write; \texttt{reset()} restores the documented reset value; and
\texttt{write\_with\_zero} (already \texttt{unsafe}) starts from all-zeroes
instead. A read-only register simply \emph{has no} \texttt{write} method---%
its marker never implements \texttt{Writable}, and the compiler rejects the
call. This is the load-bearing observation for what follows:
\textbf{conditional method availability is already the PAC's own safety
mechanism}; our change extends the same ladder by one rung.

A PAC generated by \texttt{svd2rust} defines every register method once,
generically: \texttt{write} exists for any register type implementing the
\texttt{Writable} trait. One cannot delete a method from a single
register---but one can make a method's \emph{availability conditional},
which is an idiomatic Rust pattern (a slice's \texttt{sort} exists only for
orderable elements; indeed \texttt{svd2rust} itself already rejects writes to
read-only registers by exactly this mechanism). The generator therefore
patches the crate in three small, mechanical ways:

\begin{enumerate}
\item \textbf{The shared core} (\texttt{generic.rs}, one file): the stock
  \texttt{write}, \texttt{modify}, \texttt{read}, \texttt{reset}, \ldots\
  methods gain a requirement---the register must carry an ``unconstrained''
  marker for that operation. Alongside them appear the witnessed variants
  (\texttt{write\_witnessed}, \texttt{write\_when\_ready}) and a single,
  deliberately \texttt{unsafe}, greppable escape hatch
  (\texttt{write\_unwitnessed}).
\item \textbf{Every unconstrained register}---the overwhelming
  majority---receives one-line marker implementations. Its public API is
  unchanged.
\item \textbf{Each constrained register} receives \emph{no} marker for the
  gated operation---the absence \emph{is} the gate---plus the generated
  module of \S\ref{sec:enforcement:witnesses}.
\end{enumerate}

Because the missing-marker error is a standard trait-bound failure, Rust's
diagnostic attribute turns it into a domain-specific message at the
offending line:

\begin{lstlisting}
error[E0277]: `...i2c1::cr1::CR1rs` is write-constrained by its datasheet
  --> src/i2c.rs:149
   = this register requires a witness: call `write_witnessed(f, witness)`
     or `write_when_ready(f)`; obtain the witness via `check_write_ready()`;
     bypass only with `unsafe write_unwitnessed`
\end{lstlisting}

There is no wrapper type and no method shadowing---an earlier wrapper-based
design left the stock method reachable through an ordinary reference
coercion, a silent, safe bypass; under trait gating the method \emph{does
not exist} for the constrained register, so that bypass is itself a compile
error, permanently pinned in our test suite. The patched crate is consumed
exactly like the original: it is the published, generated source plus these
additions, so a downstream project adopts it with a one-line dependency swap
and no build steps.

\paragraph{Results.}
For the running example's device, the patch touches the one shared core
file, adds one-line markers to 445 register files, and injects one
constraints module; the patched crate type-checks in ${\sim}14$\,s.
Enforcement is pinned by an eleven-program compile-fail suite: every
witness-less access fails with the message above, witness reuse and
cross-operation reuse fail with the expected ownership and type errors, and
the legal witnessed paths compile.

\subsection{Effect on Higher-Level Crates}
\label{sec:enforcement:hal}

The adoption question is whether real driver code---written by others,
before these constraints existed---can use the modified PAC \emph{as is}. We
answer it experimentally: compile the unmodified \texttt{stm32f4xx-hal} (the
community's standard hardware-abstraction layer for these chips, ${\sim}30$k
lines) against the patched PAC, consumed exactly as a downstream user would
consume it.

\paragraph{Results.}
Against the \emph{unpatched} baseline crate the HAL compiles clean---%
adoption costs nothing where nothing is constrained. Against the patched
crate there are exactly two classes of outcome and nothing else:

\begin{table}[h]
\centering
\small
\begin{tabular}{@{}llll@{}}
\toprule
class & count & where & meaning \\
\midrule
true enforcement & \textbf{14 / 14} & the HAL's I\textsuperscript{2}C driver &
  every CR1 write site fails with the datasheet \\
 & & (\texttt{i2c.rs}, \texttt{i2c/dma.rs}) &
  diagnostic; \textbf{zero false hits} elsewhere \\
\addlinespace
generic-code friction & 14 & one module (\texttt{serial/}) &
  trait-generic code must re-state its \\
 & & & requirements; ${\sim}10$ one-line bounds fix it \\
\bottomrule
\end{tabular}
\caption{Compiling the unmodified \texttt{stm32f4xx-hal} against the
constraint-injected PAC.}
\label{tab:hal-demo}
\end{table}

The first class is the running example completing its journey: the sentence
from the manual is now a compile error at the fourteen real call sites that
could violate it. Fixing such a site is a local, guided edit---the
diagnostic names the replacement:

\begin{lstlisting}
// before (no longer compiles):
self.i2c.cr1().modify(|_, w| w.stop().set_bit());
// after (checked):
self.i2c.cr1().modify_when_ready(|_, w| w.stop().set_bit())?;
\end{lstlisting}

The second class delimits ``as is'' honestly. Code that names concrete
registers---the vast majority of a HAL---needs no change unless it genuinely
violates a constraint. Code that is \emph{generic over register types} must
declare the new markers in its trait bounds, because generic code promises
its requirements up front; this HAL's serial layer is written that way and
fails to type-check even though no serial register is constrained. The cost
is small and mechanical (about ten one-line bounds in one module), but it is
a modification, and any design that makes methods conditionally available
shares it---the alternative, keeping every method callable, is precisely the
silent bypass we rejected. We therefore report adoption as two measured
numbers rather than one claim: perfect precision on concrete driver code,
and a small, quantified patch for trait-generic driver code.
